\begin{claim}{C.48}{Proved}{One boundary foliation admits a compatible intrinsic contraction}
Use the based ambient-null plane loop and the full family
$\mathscr F$ of lifted boundary fiber partitions in C.46.
For an open orbit $O$ with $|\mathscr F|=1$, there is a based
compact-open null-homotopy $K_{s,t}$ of $f_t|_O$ which
commutes with its entire original operator stabilizer and
satisfies, simultaneously for every specified boundary point,
\[
                 \Phi_zK_{s,t}=Q^{B_z}_{s,t}\Phi_z.
                                                        \tag{C.48.1}
\]
Here $Q^{B_z}$ is precisely the natural median homotopy of
C.27--C.29. The original source and target covering maps are
retained. This is an intrinsic assertion on $O$; ambient
continuity at its frontier is an additional requirement.
\end{claim}

\begin{proof}
\emph{The actual lifted operator group.}
Let $p:X\to O$ be the universal covering and choose one
real quotient $F:X\to\R$ representing the unique partition.
By C.43, $X$ is a plane, $F$ has connected proper-line fibers,
and a product chart identifies it with the second projection.
Let $Q$ be the closure, in $\operatorname{Homeo}(O)$, of the
restrictions of the original $O$-stabilizing operators.
C.30 proves that this exact closed group acts freely and
properly. Define its group of all lifts by
\[
 \widetilde Q=
 \{a\in\operatorname{Homeo}(X):pa=qp\text{ for some }q\in Q\}.
                                                        \tag{C.48.2}
\]
It includes the source deck group. Its topology is the
subspace homeomorphism-group topology, with the maps and
their inverses controlled on compact subsets.

We give the needed covering comparisons explicitly.
A lift is determined by its value at one point once its
downstairs map is specified. For a convergent sequence of
downstairs maps and a convergent compatible choice of those
values, the lifts converge on compact subsets: cover a
compact path from that point by finitely many covering
charts and continue the chosen sheet successively; the
same argument on finitely many such neighborhoods controls
a compact set. Apply it to the inverse maps as well.
Conversely a convergent sequence of lifts gives convergence
downstairs on each compact set by finitely many local
sections of $p$. These observations prove that (C.48.2) is
closed and that its projection to $Q$ has the stated
topology of lifts.

The action of $\widetilde Q$ on $X$ is free. A lift fixing
a point has downstairs map fixing a point, hence equal to
the identity; the resulting deck map fixing a point is the
identity too. It is proper in this topology. Indeed, from
$x_n\to x$ and $a_nx_n\to y$, properness downstairs gives
a subsequence $q_n\to q$. There is a unique lift $a$ of $q$
with $ax=y$, and the covering-chart argument gives
$a_n\to a$. The same argument makes evaluation at a point
a homeomorphism onto a closed subset of $X$: convergence
of its values supplies convergence of the corresponding
downstairs maps and then of the selected lifts.
Thus $\widetilde Q$ is locally compact. Its faithful product
action on $X\times\R$ makes it a Lie group by the already
checked Hilbert--Smith theorem of C.33. Its dimension is
at most two, since each orbit is properly embedded in $X$.

The original operators preserve the unique partition by
(C.46.5). This property passes to their restrictions'
closure and to all lifts. More explicitly, choose lifts
of an approximating sequence whose values at one point
converge to the chosen lift. Equality of $F$-values on a
fiber passes to the limit, and the inverse maps give
equality of the resulting fiber partitions. Since $F$ is
an open quotient, every $a\in\widetilde Q$ therefore has a
unique induced line homeomorphism $\bar a$ with
\[
                         Fa=\bar aF.                   \tag{C.48.3}
\]
A continuous section of $F$ gives $\bar a=Fa\sigma$, and
the inverse formula uses $a^{-1}$; hence this induced
homomorphism is continuous in the actual topologies.

Lift the original restricted isotopy from the identity
to $a_t$ on $X$. C.32 proves its null evaluation loop in
$O$, so $a_1=\id$. Every element of (C.48.2) commutes with
$a_t$: its commutator is a continuous deck-valued function
of $t$, initially the identity. Write
\[
                         Fa_t=b_tF.                    \tag{C.48.4}
\]
These are the based target lifts, and $b_0=b_1=\id$ by
C.29 for a circle target and directly for a line target.
The group generated by the paths $b_t$ acts transitively
on $\R$. Its orbits are intervals, because every word is
joined to the identity by a path of words. A finite
endpoint of a proper orbit interval would be fixed by
every word. Such a common fixed point would project to a
fixed point of the nontrivial transitive boundary orbit,
which is impossible. Thus the orbit interval is all of
$\R$.

\emph{The base operator image and its kernel.}
Take the closed transitive group generated on this real
line. The evaluation proof of C.30 applies to the line
without change and shows that its full centralizer is
locally compact, free and proper. Its product action on
$\R^3$ makes it Lie. No decreasing element is possible:
it would have a fixed point, contrary to freeness.
In dimension one its identity component acts transitively
on the line, and freeness excludes additional components;
it is the real translation group. In dimension zero the
nearest-orbit-point argument of C.44 gives either the
identity group or a cyclic translation group. Choose the
resulting line coordinate. We have thereby proved, with
the comparison (C.48.3), that every base operator is a
translation. Write
\[
 Fa=F+\alpha(a),\qquad
 \alpha:\widetilde Q\longrightarrow(\R,+),\qquad
 K=\ker\alpha,
                                                        \tag{C.48.5}
\]
after this specified change of target coordinate.

Restriction of $K$ to any fiber is faithful, free and
proper. Evaluation gives it the same topology as its
closed orbit in that line. The preceding one-dimensional
classification therefore gives exactly
\[
                         K=1,\quad K\cong\Z,
                         \quad\text{or}\quad K\cong\R.
                                                        \tag{C.48.6}
\]
All coordinate changes below are homeomorphisms of the
source and reparametrizations of this same target.
We construct them and keep the action of every component.

\emph{The real kernel.}
For $K\cong\R$, its orbits are the entire fibers. A global
section of its principal bundle over the line, constructed
as in C.33, gives coordinates $(u,r)$ with $F(u,r)=r$ and
\[
                         k_c(u,r)=(u+c,r).              \tag{C.48.7}
\]
The quotient $\widetilde Q/K$ acts freely and properly on
$X/K=\R$, by the exact local-section argument in C.33.
Its image under $\alpha$ is consequently closed, and is
$0$, $d\Z$ with $d>0$, or $\R$.
For image zero (C.48.7) is the complete operator action.
For image $d\Z$, choose a label-$d$ element. Its action is
\[
                 q(u,r)=(\lambda u+\tau(r),r+d),
                         \lambda\in\R^\times.          \tag{C.48.8}
\]
The scalar is constant because conjugation on the actual
group $K$ is a single continuous group automorphism.
Choose $c$ on $[0,d]$ with
$c(d)=\lambda c(0)+\tau(0)$ and extend it by
$c(r+d)=\lambda c(r)+\tau(r)$. In the coordinates
$u'=u-c(r)$ the complete additional operator is
\[
                         q(u',r)=(\lambda u',r+d).
                                                        \tag{C.48.9}
\]
This retains negative $\lambda$ as well.

For image $\R$, the group $\widetilde Q$ acts transitively
and freely on $X$. Both its kernel and quotient are
connected lines, so it is connected. A Lie-algebra element
mapping to $1$ integrates to a section $s_r$ with
$\alpha(s_r)=r$. Conjugation on $K$ has the form
$s_rk_us_r^{-1}=k_{e^{cr}u}$ for a real constant $c$.
Evaluation gives the exact group coordinates and product
\[
 X\cong\widetilde Q\cong\R^2,\qquad
 (u,r)(v,z)=(u+e^{cr}v,r+z),\qquad F(u,r)=r.
                                                        \tag{C.48.10}
\]
The formula proves directly that these coordinates are
global; $k_us_r$ is the unique element with the specified
base and fiber values. The operator action is left
multiplication, and its full centralizer consists of
right multiplications.

\emph{The cyclic kernel.}
For $K\cong\Z$, the construction (C.44.5) supplies
coordinates with $F(u,r)=r$ and generator
$k(u,r)=(u+1,r)$. The quotient is the actual cylinder
$X/K=(\R/\Z)\times\R$. The action of
$\widetilde Q/K$ on it is free and proper: lift convergent
point sequences by local sections and correct by elements
of $K$, then apply properness upstairs, exactly as in
C.33. The compact circle fibers show that its base image
is closed. Indeed, a sequence of base translations in a
compact interval sends one fixed point into a compact
cylinder strip; properness gives a convergent subsequence
of the quotient elements. Thus again the image is
$0$, $d\Z$, or $\R$.

For image $d\Z$, the additional generator has the form
\[
 q(u,r)=(g_r(u),r+d),\qquad
 g_r(u+1)=g_r(u)+\epsilon,\quad\epsilon\in\{1,-1\}.
                                                        \tag{C.48.11}
\]
Here $g_r$ is increasing for $\epsilon=1$ and decreasing
for $\epsilon=-1$. Choose an increasing unit-translation
equivariant path $\rho_r$ on $0\leq r\leq d$ from
$\rho_0=\id$ to $\rho_d=g_0\circ(\epsilon\id)$.
The convex interpolation of these two increasing maps
is such a path. Extend it in both directions by
\[
                 \rho_{r+d}=g_r\rho_r(\epsilon\id).
                                                        \tag{C.48.12}
\]
The endpoint equation makes this extension continuous.
Every $\rho_r$ remains increasing and commutes with unit
translation. The homeomorphism $(u,r)\mapsto(\rho_r(u),r)$
therefore changes the two generators into exactly
\[
                 k(u,r)=(u+1,r),\qquad
                 q(u,r)=(\epsilon u,r+d).              \tag{C.48.13}
\]

For base image $\R$, the projection with discrete kernel
is a covering homomorphism of Lie groups. Its identity
component maps isomorphically onto the simply connected
line. It centralizes $K$, since its conjugation map to
$\operatorname{Aut}(\Z)=\{1,-1\}$ is continuous and is
initially the identity. Thus
$\widetilde Q=\Z\times\R$. Use this real subgroup to
transport the fiber $F^{-1}(0)$ to every other fiber,
and use the nearest-orbit-point coordinate for $K$ on
that first fiber. These are global coordinates in which
\[
           k(u,r)=(u+1,r),\qquad s_z(u,r)=(u,r+z).
                                                        \tag{C.48.14}
\]
For base image zero only $k$ remains.

\emph{Trivial kernel and the possible dense base image.}
For $K=1$, the group $\widetilde Q$ embeds algebraically
in the translation group and is abelian. A positive
dimensional identity component maps isomorphically onto
$\R$; injectivity excludes other components. The same
transport from $F^{-1}(0)$ gives just $s_z$ in (C.48.14).
In dimension zero, $X\to X/\widetilde Q$ is a regular
surface covering with simply connected source. Its deck
group is the actual fundamental group of the quotient.
A noncompact surface has free fundamental group, and
an abelian free group is trivial or infinite cyclic.
For a compact quotient the surface classification gives
only the torus with abelian torsion-free fundamental
group: the sphere and projective plane do not have plane
universal cover, the Klein-bottle group is nonabelian,
and every remaining closed surface group is nonabelian.
Consequently the discrete possibilities are exactly
$1$, $\Z$, and $\Z^2$.

The cyclic case has a nonzero base translation; choose
its positive generator. The construction
(C.48.11)--(C.48.13), now without the unit-translation
condition, changes it to $(u,r)\mapsto(\epsilon u,r+d)$.
The trivial case has no equations.
For $\widetilde Q\cong\Z^2$ the quotient is a torus,
and its base image is dense: an injective homomorphism
from $\Z^2$ into $\R$ has two generators with irrational
ratio. The foliation descends to an oriented minimal
topological foliation of that torus. To verify minimality,
the saturation of any lifted fiber under the deck group
is $F^{-1}(r+\alpha(\widetilde Q))$, which is dense because
$F$ is open and the displayed set of levels is dense.
The quotient torus is orientable, and its deck maps
preserve both the surface and the transverse orientation,
so the leaf orientation descends as asserted.

We give the elementary foliation conjugacy needed here,
including its maps. It is the minimal-torus conclusion
recorded in \cite[Proposition~2.5]{C_MionMouton2025}; no
bi-foliation classification or asymptotic-cycle statement
is needed. The oriented local flowboxes give a nonsingular
continuous flow with these leaves by the construction
used in C.43. Choose a small closed transverse interval,
with endpoints off the orbit of an interior point, and a
first-return arc from that point to a different point of
the interval. A flowbox excludes sufficiently short
positive returns; minimality supplies a later return,
and compactness gives a first one. The endpoints chosen
off the orbit put that return in the interval's interior.
The arc's interior avoids the interval, and its endpoint
differs from its initial point because there is no
periodic orbit. Its union with the intervening
transverse subinterval is a simple closed curve.
In a thin flow strip along the arc, tilt the arc strictly
in the transverse direction pointing from the return
point back to the initial point. The tilt may be chosen
arbitrarily small, because the two endpoint transverse
coordinates are distinct; choose it so that the new
terminal point remains strictly between the old
endpoints in the return interval. Join the new endpoints
inside the transverse interval. Both pieces now have
the same strict transverse direction; straightening
their two corners in flowboxes gives an embedded closed
transversal $\gamma$.

Every leaf meets $\gamma$, since it is dense and a
flowbox around $\gamma$ is open. Forward and backward
returns exist: on a compact minimal flow each forward
and backward limit set is a nonempty closed invariant
set, hence the entire torus. Flowboxes and avoidance
of $\gamma$ on a compact interior return arc show that
the first return time is continuous and positive.
The return map $h:\gamma\to\gamma$ is a minimal
orientation-preserving circle homeomorphism. The map
which follows each point for a fraction of its return
time gives the leaf-preserving homeomorphism from its
mapping torus
\[
 [0,1]\times\gamma/((1,z)\sim(0,hz))\longrightarrow X/\widetilde Q.
                                                        \tag{C.48.15}
\]
Surjectivity uses the preceding returns; uniqueness of
the last intersection gives injectivity, and the compact
source and Hausdorff target give the inverse continuity.

For completeness, choose a probability measure invariant
under $h$ by taking a weak limit of finite orbit averages.
Minimality gives full support, and the absence of periodic
points gives no atoms. Its cumulative distribution is
therefore a circle homeomorphism. It conjugates $h$ to a
Lebesgue-measure-preserving increasing homeomorphism,
hence to a rotation: preservation of the measure of
every interval makes the displacement constant.
The angle $\theta$ is irrational by minimality.
After this conjugacy, the explicit map
\[
              [s,z]\longmapsto[(s,z+s\theta)]
                         \quad\text{in }\R^2/\Z^2
                                                        \tag{C.48.16}
\]
is well defined at the identified endpoints and sends
the suspended leaves to the lines of slope $\theta$.
This proves the required foliation conjugacy.

Lift (C.48.15)--(C.48.16) to the universal covers. It
intertwines the original deck action with translations
by $\Z^2$, through the induced isomorphism of their
fundamental groups. Each original $F$-fiber becomes a
whole affine line of slope $\theta$. A linear transverse
coordinate $r$ and a complementary linear coordinate
$u$ thus give
\[
       q(u,r)=(u+\beta(q),r+\alpha'(q))\quad(q\in\widetilde Q),
                                                        \tag{C.48.17}
\]
where these translation vectors form a lattice, and
$\alpha'(\widetilde Q)$ is dense. The original function
$F$ is a line homeomorphism composed with $r$: both
are open quotients with exactly the same connected
fibers. This gives the exact target reparametrization,
rather than identifying the two quotient functions
without a comparison map.

\emph{Contraction in the established normal forms.}
Except for the separately displayed transitive case
(C.48.10), write the actual lifted loop as
\[
                 a_t(u,r)=(h_t(u,r),b_t(r)).             \tag{C.48.18}
\]
Each $h_t(\,\cdot\,,r)$ is an increasing line
homeomorphism. Its endpoints are the identity family.
Let $B_{s,t}$ be the natural median contraction of
$b_t$ from C.27, and set
\[
 H_{s,t}(u,r)=
       \bigl((1-s)h_t(u,r)+su,\ B_{s,t}(r)\bigr).
                                                        \tag{C.48.19}
\]
We check every operator equation rather than changing
the operator group. For $k(u,r)=(u+1,r)$ the equation
is $h_t(u+1,r)=h_t(u,r)+1$, which (C.48.19) preserves.
For $q(u,r)=(\epsilon u,r+d)$ it is
\[
 h_t(\epsilon u,r+d)=\epsilon h_t(u,r),\qquad
 b_t(r+d)=b_t(r)+d.
                                                        \tag{C.48.20}
\]
Convex interpolation preserves the first identity and
median naturality preserves the second.
For the real kernel (C.48.7), commutation says
$h_t(u,r)=u+\phi_t(r)$; with (C.48.9) the further
equation is $\phi_t(r+d)=\lambda\phi_t(r)$.
Scaling $\phi_t$ by $1-s$ preserves it, including
negative $\lambda$.
For the real base translations in (C.48.14), the
equations are that $h_t$ is independent of $r$ and
$b_t(r)=r+c_t$. Both are preserved by (C.48.19).

For the lattice (C.48.17), the dense set of base
translations and continuity force $b_t(r)=r+c_t$.
The fiber equation is exactly
\[
 h_t(u+\beta(q),r+\alpha'(q))=h_t(u,r)+\beta(q).
                                                        \tag{C.48.21}
\]
It is preserved by convex interpolation with $u$.
The median of translations is again a translation,
so the base equation is preserved as well.
The trivial operator case has no extra equation.
These cases exhaust (C.48.6).

In (C.48.10), write the loop as right multiplication
by $(v_t,c_t)$. Its action is
$(u,r)\mapsto(u+e^{cr}v_t,r+c_t)$.
The median base contraction is translation by some
$C_{s,t}$. Right multiplication by
\[
                         ((1-s)v_t,C_{s,t})             \tag{C.48.22}
\]
is the required contraction and still commutes with
every left multiplication.

Every map in (C.48.19) is a homeomorphism: its base
map is a line homeomorphism, and its fiber map is
an increasing convex combination of two increasing
onto maps. The maps are jointly continuous. Their
inverses are jointly continuous as well: on compact
parameter sets, monotonicity and two fixed bracketing
fiber values bound the inverse, and continuity and
uniqueness identify every limit. This is the inverse
continuity argument of C.31 applied to the displayed
fiber family and the continuous inverse base family.
Formula (C.48.22) has the same properties directly
from the group law. All endpoint identities are
preserved. We have therefore obtained a based
null-homotopy commuting with all of $\widetilde Q$.

In particular it commutes with the original source
deck group, so it descends to a compact-open
null-homotopy $K_{s,t}$ on $O$, with continuous
inverse family. Its commutation with every original
operator follows from (C.48.2), retaining the
restriction kernel, which acts identically on $O$.
Finally each other boundary real quotient is an
exact reparametrization of $F$ by the first case
of C.45. Median naturality under that line
homeomorphism gives its required quotient identity.
For a circle target it commutes with the target
integer translations and descends. The equations
$q_zF_z=\Phi_zp$ now give (C.48.1) for every $z$.
No limit across the ambient frontier has been
deduced from this intrinsic construction.
\end{proof}
